\section{术语消化：本期关键词索引}

前面的章节横跨技术、制度、组织和资本市场。这里把关键词集中消化，方便后续和大模型季报、Agent 综述、产品访谈互相索引。

\begin{longtable}{p{0.24\textwidth}p{0.32\textwidth}p{0.35\textwidth}}
\toprule
术语 & 一句话解释 & 在本期中的作用 \\
\midrule
认知智能 & 感知智能之后的理解、推理和知识阶段 & 智谱早期方向锚点。 \\
P2P & Paper to Project / Product & 清华成果转化传统。 \\
科技成果转化 & 把研究成果合规转成公司资产和产品 & 智谱成立的制度路径。 \\
GPT-3 & 规模范式的早期外部冲击 & 改变大模型技术想象。 \\
ChatGPT & 大模型用户心智爆发节点 & 让产业节奏突然加速。 \\
Scaling Law & 模型能力随规模和反馈扩展的规律 & 大模型竞争的底层范式。 \\
Post-training & 预训练之后用偏好、任务、工具和反馈继续训练 & 模型从“会说”走向“可用”的关键阶段。 \\
RL & Reinforcement Learning，强化学习 & 张鹏认为会带来新的计算方式和能力增长轴。 \\
Agent & 调用工具并执行任务的智能体 & 连接模型能力与真实应用。 \\
DeepSeek & 代表开源、效率和工程冲击的行业样本 & 让模型公司重新理解发布和开源。 \\
开源/闭源 & 生态扩散与商业控制的策略组合 & 模型公司必须取舍。 \\
Benchmark & 用统一任务比较模型能力的评测 & 有助比较，但不能替代真实使用。 \\
真实使用 & 开发者、客户和产品对模型的持续采用 & 上市公司阶段更关键的能力证明。 \\
IPO & 上市，进入公开市场和更强合规治理 & 智谱的新阶段。 \\
先行者 & 早投入、早铺路、早承担不确定性的人 & 张鹏希望智谱留下的历史注脚。 \\
\bottomrule
\end{longtable}

\subsection{本章小结}

这些关键词共同描述了一家模型公司的完整链路：从实验室、成果转化、技术范式冲击，到开源生态、商业化、上市治理和长期 AGI 叙事。

